\documentclass[12pt]{article}
\usepackage[utf8]{inputenc}
\usepackage[T1]{fontenc}
\usepackage{hyperref}
\title{Kreye Language Resources for K-12 Curriculum Development}
\author{Compiled by Claude Sonnet 5.5 (claude-sonnet-5-5) for Andr\'as Kornai}
\date{}
\begin{document}
\maketitle

\section{Introduction}
Kreye (Kr\~eje, Kr\'ey\'e; ISO 639-3 xre) is a Je language that was spoken in Maranh\~ao and Par\'a, Brazil. Sources disagree on its status. Wikipedia lists it as moribund, with about 15 speakers in 2012, while ASJP records it as extinct since 2010. Joshua Project lists a people group of about 100 people who mainly speak Portuguese. This report rests on a single search round, and documentation appears very thin.

\section{Primary Educational Resources}
None were found.

\section{Linguistic and Reference Materials}
\begin{itemize}
\item Wikipedia, Kr\~eje language: \url{https://en.wikipedia.org/wiki/Kr%E1%BA%BDje_language}
\item Curt Nimuend\'aj\'u's vocabulary (published 1914 in \textit{Zeitschrift f\"ur Ethnologie}) is the earliest substantial record. A Unicamp conference abstract describes work on making it usable for comparative research: \url{https://www.prp.unicamp.br/pibic/congressos/xxicongresso/resumos/104778.pdf}
\item Numerals entry, Max Planck EVA: \url{https://lingweb.eva.mpg.de/channumerals/Kreye.htm}
\item ASJP word list entry: \url{https://asjp.clld.org/languages/KREYE}
\end{itemize}

\section{Religious and Community Resources}
\begin{itemize}
\item Joshua Project profile of the Kreje people in Brazil: \url{https://joshuaproject.net/people_groups/11240}
\end{itemize}

\section{Digital and Audio Resources}
No audio recordings or digital teaching resources were found.

\section{Curriculum Development Context}
The language is documented mainly by a single early vocabulary and short word lists. A heritage curriculum would have to be built around cultural content and basic vocabulary, together with descendant community members and the Brazilian Je-language research community. It could not rest on a reference grammar or on recorded material, because the search found neither. Claims about the language's content should be checked against the primary sources.

\end{document}